\subsection{Критерий принадлежности числа спектру самосопряженного оператора. Вещественность спектра самосопряженного оператора}

\begin{definition}[Оператор \(A_\lambda\)]
	Пусть \(H\) --- гильбертово пространство над \(\mathbb{C}\),
	\[
	A\in\mathcal{L}(H),
	\qquad
	\lambda\in\mathbb{C}.
	\]
	Обозначим
	\[
	A_\lambda=A-\lambda I.
	\]
\end{definition}

\begin{remark}[Используемые факты]
	В доказательстве используем два уже доказанных факта.
	
	\begin{itemize}
		\item Если оператор ограничен снизу, то он инъективен и его образ замкнут.
		
		\item Для самосопряженного оператора \(A\) выполнено разложение
		\[
		H=\overline{\operatorname{Im}A_\lambda}\oplus_{\perp}\ker A_\lambda.
		\]
	\end{itemize}
\end{remark}

\begin{theorem}[Критерий принадлежности числа спектру]
	Пусть \(H\) --- гильбертово пространство над \(\mathbb{C}\), и
	\[
	A\in\mathcal{L}(H)
	\]
	--- самосопряженный оператор. Тогда:
	
	\begin{itemize}
		\item \[
		\lambda\in\rho(A)
		\Longleftrightarrow
		A_\lambda \text{ ограничен снизу},
		\]
		то есть
		\[
		\exists m>0
		\quad
		\forall x\in H
		\qquad
		\|A_\lambda x\|\geqslant m\|x\|.
		\]
		
		\item \[
		\lambda\in\sigma(A)
		\Longleftrightarrow
		\exists \{x_n\}\subset H:
		\|x_n\|=1,
		\quad
		A_\lambda x_n\to0.
		\]
	\end{itemize}
\end{theorem}

\begin{proof}
	Сначала докажем критерий через ограниченность снизу, а затем аккуратно перепишем его
	отрицание.
	
	\textbf{1. Если \(\lambda\in\rho(A)\), то \(A_\lambda\) ограничен снизу.}
	
	Пусть
	\[
	\lambda\in\rho(A).
	\]
	Тогда по определению резольвентного множества существует ограниченный обратный оператор
	\[
	A_\lambda^{-1}\in\mathcal{L}(H).
	\]
	По критерию из билета 8.1 отсюда следует, что \(A_\lambda\) ограничен снизу.
	
	То есть существует \(m>0\), такое что
	\[
	\|A_\lambda x\|\geqslant m\|x\|
	\qquad
	\forall x\in H.
	\]
	
	\textbf{2. Если \(A_\lambda\) ограничен снизу, то \(\lambda\in\rho(A)\).}
	
	Пусть существует \(m>0\), такое что
	\[
	\|A_\lambda x\|\geqslant m\|x\|
	\qquad
	\forall x\in H.
	\]
	Нужно доказать, что
	\[
	A_\lambda^{-1}\in\mathcal{L}(H).
	\]
	По билету 8.1 для этого достаточно показать, что
	\[
	\operatorname{Im}A_\lambda=H.
	\]
	
	\begin{itemize}
		\item \textbf{Ядро тривиально.}
		
		Если
		\[
		x\in\ker A_\lambda,
		\]
		то
		\[
		A_\lambda x=0.
		\]
		По ограниченности снизу
		\[
		0=\|A_\lambda x\|\geqslant m\|x\|.
		\]
		Так как \(m>0\), получаем
		\[
		x=0.
		\]
		Значит,
		\[
		\ker A_\lambda=\{0\}.
		\]
		
		\item \textbf{Образ плотен.}
		
		По разложению из билета 11.2
		\[
		H=\overline{\operatorname{Im}A_\lambda}\oplus_{\perp}\ker A_\lambda.
		\]
		Так как
		\[
		\ker A_\lambda=\{0\},
		\]
		получаем
		\[
		H=\overline{\operatorname{Im}A_\lambda}.
		\]
		
		\item \textbf{Образ замкнут.}
		
		Так как \(A_\lambda\) ограничен снизу, то по следствию из билета 8.1 его образ замкнут:
		\[
		\operatorname{Im}A_\lambda
		\text{ замкнут}.
		\]
		Но уже известно, что
		\[
		\overline{\operatorname{Im}A_\lambda}=H.
		\]
		Следовательно,
		\[
		\operatorname{Im}A_\lambda=H.
		\]
	\end{itemize}
	
	Итак, \(A_\lambda\) инъективен и сюръективен на \(H\), а также ограничен снизу. Поэтому по
	критерию из билета 8.1
	\[
	A_\lambda^{-1}\in\mathcal{L}(H).
	\]
	Значит,
	\[
	\lambda\in\rho(A).
	\]
	
	\textbf{3. Перепишем отрицание критерия.}
	
	Из первого пункта имеем
	\[
	\lambda\in\sigma(A)
	\Longleftrightarrow
	\lambda\notin\rho(A)
	\Longleftrightarrow
	A_\lambda \text{ не ограничен снизу}.
	\]
	Отрицание ограниченности снизу имеет вид:
	\[
	\forall m>0
	\quad
	\exists x\in H
	\qquad
	\|A_\lambda x\|<m\|x\|.
	\]
	Такой \(x\) автоматически ненулевой, иначе получилось бы \(0<0\).
	
	Выберем
	\[
	m=\frac{1}{n}.
	\]
	Тогда существует \(u_n\neq0\), такое что
	\[
	\|A_\lambda u_n\|<\frac{1}{n}\|u_n\|.
	\]
	Положим
	\[
	x_n=\frac{u_n}{\|u_n\|}.
	\]
	Тогда
	\[
	\|x_n\|=1
	\]
	и
	\[
	\|A_\lambda x_n\|
	=
	\frac{\|A_\lambda u_n\|}{\|u_n\|}
	<
	\frac{1}{n}.
	\]
	Следовательно,
	\[
	A_\lambda x_n\to0.
	\]
	
	Обратно, если существует последовательность \(\{x_n\}\), такая что
	\[
	\|x_n\|=1,
	\qquad
	A_\lambda x_n\to0,
	\]
	то \(A_\lambda\) не может быть ограничен снизу. Действительно, если бы существовало
	\(m>0\), такое что
	\[
	\|A_\lambda x\|\geqslant m\|x\|
	\qquad
	\forall x\in H,
	\]
	то для \(x=x_n\) получили бы
	\[
	\|A_\lambda x_n\|\geqslant m,
	\]
	что противоречит сходимости
	\[
	A_\lambda x_n\to0.
	\]
	Значит,
	\[
	\lambda\in\sigma(A).
	\]
\end{proof}

\begin{theorem}[Вещественность спектра самосопряженного оператора]
	Пусть \(H\) --- гильбертово пространство над \(\mathbb{C}\), и
	\[
	A\in\mathcal{L}(H)
	\]
	--- самосопряженный оператор. Тогда
	\[
	\sigma(A)\subset\mathbb{R}.
	\]
\end{theorem}

\begin{proof}
	Докажем равносильное утверждение:
	\[
	\lambda\notin\mathbb{R}
	\Longrightarrow
	\lambda\in\rho(A).
	\]
	
	\textbf{1. Представим \(\lambda\) и выделим самосопряженную часть.}
	
	Пусть
	\[
	\lambda=\mu+i\nu,
	\qquad
	\mu,\nu\in\mathbb{R},
	\qquad
	\nu\neq0.
	\]
	Тогда
	\[
	A_\lambda
	=
	A-\lambda I
	=
	A-\mu I-i\nu I.
	\]
	Обозначим
	\[
	A_\mu=A-\mu I.
	\]
	Так как \(A\) самосопряжен и \(\mu\in\mathbb{R}\), оператор \(A_\mu\) тоже самосопряжен.
	
	\textbf{2. Докажем оценку снизу для \(A_\lambda\).}
	
	Возьмём произвольный \(x\in H\). Тогда
	\[
	A_\lambda x
	=
	A_\mu x-i\nu x.
	\]
	Поэтому
	\[
	\|A_\lambda x\|^2
	=
	(A_\mu x-i\nu x,\ A_\mu x-i\nu x).
	\]
	Раскрывая скалярное произведение, получаем
	\[
	\|A_\lambda x\|^2
	=
	\|A_\mu x\|^2
	-i\nu(x,A_\mu x)
	+i\nu(A_\mu x,x)
	+|\nu|^2\|x\|^2.
	\]
	Так как \(A_\mu\) самосопряжен, по билету 11.1
	\[
	(A_\mu x,x)\in\mathbb{R}.
	\]
	Значит,
	\[
	(x,A_\mu x)=(A_\mu x,x),
	\]
	и смешанные слагаемые сокращаются:
	\[
	-i\nu(x,A_\mu x)
	+i\nu(A_\mu x,x)=0.
	\]
	Следовательно,
	\[
	\|A_\lambda x\|^2
	=
	\|A_\mu x\|^2
	+
	|\nu|^2\|x\|^2.
	\]
	Отсюда
	\[
	\|A_\lambda x\|^2
	\geqslant
	|\nu|^2\|x\|^2.
	\]
	Значит,
	\[
	\|A_\lambda x\|
	\geqslant
	|\nu|\|x\|.
	\]
	То есть \(A_\lambda\) ограничен снизу.
	
	\textbf{3. Применим критерий принадлежности спектру.}
	
	По доказанному критерию
	\[
	A_\lambda \text{ ограничен снизу}
	\Longrightarrow
	\lambda\in\rho(A).
	\]
	Следовательно,
	\[
	\lambda\notin\mathbb{R}
	\Longrightarrow
	\lambda\in\rho(A).
	\]
	То есть нерегулярных точек вне вещественной оси нет, поэтому
	\[
	\sigma(A)\subset\mathbb{R}.
	\]
\end{proof}
